\documentclass[11pt, aspectratio=169]{beamer}

\usetheme{metropolis}
\metroset{
  numbering=fraction,
  progressbar=foot,
  block=fill,
  sectionpage=none,
}

\usepackage{fontspec}
\usepackage{polyglossia}
\setdefaultlanguage{russian}
\setotherlanguage{english}

\setmainfont{Arial}
\setsansfont{Arial}
\setmonofont{Fira Code}

\usepackage{graphicx}
\usepackage{qrcode}

\usepackage{tikz}
\usetikzlibrary{positioning, arrows.meta}

\definecolor{acceptedGreen}{HTML}{2E7D32}
\definecolor{acceptedFill}{HTML}{E8F5E9}
\definecolor{rejectedGray}{HTML}{757575}
\definecolor{rejectedFill}{HTML}{F5F5F5}
\definecolor{taskBlue}{HTML}{1565C0}
\definecolor{taskFill}{HTML}{E3F2FD}
\definecolor{gapRed}{HTML}{C62828}
\definecolor{slackPurple}{HTML}{4A154B}
\definecolor{mutedText}{HTML}{616161}

\definecolor{slackCardBorder}{HTML}{DDDDDD}

\hypersetup{
  colorlinks=true,
  urlcolor=taskBlue,
  linkcolor=.,
}
\setbeamerfont{footnote}{size=\scriptsize}

\tikzset{
  artifact accepted/.style={
    draw=acceptedGreen, line width=0.8pt,
    fill=acceptedFill,
    rounded corners=2pt,
    align=center, font=\small\sffamily,
    inner sep=4pt
  },
  artifact rejected/.style={
    draw=rejectedGray, line width=0.6pt, dashed,
    fill=rejectedFill,
    rounded corners=2pt,
    align=center, font=\small\sffamily, text=rejectedGray,
    inner sep=4pt
  },
  artifact task/.style={
    draw=taskBlue, line width=0.6pt,
    fill=taskFill,
    rounded corners=2pt,
    align=center, font=\footnotesize\sffamily,
    inner sep=3pt
  },
  fk arrow/.style={
    -{Stealth[length=2mm, width=2mm]},
    line width=0.7pt,
    draw=black!60
  },
  gap box/.style={
    draw=gapRed, line width=0.8pt, dashed,
    fill=gapRed!5,
    rounded corners=3pt,
    align=center,
    inner sep=8pt
  }
}

\newcommand{\artid}[1]{\texttt{\footnotesize #1}}
\newcommand{\shot}[1]{%
  \begin{center}
    \includegraphics[width=0.92\textwidth, height=0.74\textheight, keepaspectratio]{img/shots/#1}
  \end{center}
}

\title{Intent Debt}
\subtitle{Потерянная цель и потерянный код}
\author{\textit{Максим Смирнов}}
\date{}

\begin{document}

\begin{frame}[plain]
  \vfill
  \begin{center}
    {\Huge\textbf{Intent Debt}}\\[0.35cm]
    {\large\textcolor{mutedText}{Потерянная цель и потерянный код}}\\[1.1cm]
    {\normalsize\textit{Максим Смирнов}}
  \end{center}
  \vfill
\end{frame}

% ---------------------------------------------------------------- 2
\begin{frame}{Месяц назад}
  % A real Mattermost window (local instance), not a mock-up: two aligned
  % crops — channel header and conversation tail. Re-shoot: see mm/.
  \begin{center}
  \setlength{\fboxsep}{0pt}\setlength{\fboxrule}{0.6pt}%
  \fcolorbox{slackCardBorder}{white}{%
    \begin{minipage}{0.76\textwidth}
      \includegraphics[width=\linewidth]{img/mattermost-header.png}\par\nointerlineskip
      \includegraphics[width=\linewidth]{img/mattermost-message.png}%
    \end{minipage}}
  \end{center}

  \begin{center}
  \small
  \textcolor{mutedText}{\textit{Через шесть недель фича в проде. Через полгода на ретро:
  <<а precision на thin-file вырос?>>}}\\
  \textcolor{mutedText}{\textit{Ответить некому --- нигде не записано, что мы этого ждали.}}
  \end{center}
\end{frame}

% ---------------------------------------------------------------- 3
\begin{frame}{Что теряется по дороге}
  \vspace{0.2cm}
  \begin{center}
  \begin{tikzpicture}
    \node[draw=slackPurple, line width=0.6pt, fill=slackPurple!10,
          rounded corners=2pt, minimum width=3cm, minimum height=1.1cm,
          align=center, font=\small]
          (start) at (0, 0) {Разговор:\\<<сделаем X, чтобы стало Y>>};

    \node[draw=taskBlue, line width=0.6pt, fill=taskFill,
          rounded corners=2pt, minimum width=3cm, minimum height=1.1cm,
          align=center, font=\small]
          (end) at (8, 0) {Результат: \\код в проде};

    \draw[-{Stealth[length=4mm, width=4mm]}, line width=1.5pt, black!50]
      (start.east) -- (end.west);

    \node[gap box, minimum width=5.5cm] (gap) at (4, -2.2) {%
      \textcolor{gapRed}{\textbf{пропасть}}\\[4pt]
      \footnotesize какую метрику двигаем? какая была гипотеза?\\
      \footnotesize что отбросили по дороге --- и почему?
    };
    \draw[-{Stealth[length=2mm]}, dashed, draw=gapRed, line width=0.7pt]
      (gap.north) -- (4, -0.2);
  \end{tikzpicture}
  \end{center}
\end{frame}

% ---------------------------------------------------------------- 4
\begin{frame}{У этой пропасти есть имя}
  \vspace{0.2cm}

  \begin{center}
  {\Huge \textbf{Intent Debt}}
  \end{center}

  \vspace{0.3cm}

  Термин предложила Маргарет-Энн Стори: рядом с техническим и когнитивным долгом
  она ставит третий\footnote{M.-A.\ Storey. \textit{From Technical Debt to Cognitive
  and Intent Debt: Rethinking Software Health in the Age of AI.} ACM Queue, 2026.
  \href{https://doi.org/10.1145/3807966}{\texttt{doi.org/10.1145/3807966}}
  (препринт: \href{https://arxiv.org/abs/2603.22106}{\texttt{arXiv:2603.22106}}).} ---

  \vspace{0.3cm}

  \begin{center}
  \begin{minipage}{0.85\textwidth}
  \small
  \textit{<<the absence or erosion of explicit rationale, goals,
  and constraints that guide how a system evolves>>}
  \end{minipage}
  \end{center}

  \vspace{0.35cm}

  \begin{center}
  \footnotesize\textcolor{mutedText}{%
    Технический долг --- это про то, \emph{как} написан код.\\
    Intent debt --- про то, \emph{зачем} он вообще написан.
  }
  \end{center}
\end{frame}

% ---------------------------------------------------------------- 5
\begin{frame}{Потерянный код}
  \vspace{0.1cm}

  Ещё в 1985 году Питер Наур писал\footnote{P.\ Naur. \textit{Programming as Theory
  Building.} 1985.
  \href{https://pages.cs.wisc.edu/~remzi/Naur.pdf}{\texttt{pages.cs.wisc.edu/\textasciitilde remzi/Naur.pdf}}}:
  программирование --- это построение \emph{теории} того,
  как задача решается программой.

  \vspace{0.3cm}

  \begin{itemize}
    \setlength\itemsep{0.35em}
    \item Живая система --- это теория в головах команды. Код --- только её проекция.
    \item Когда команда уходит, теория умирает вместе с ней. Код остаётся.
    \item Такой код компилируется, тесты зелёные --- и всё равно он \textbf{потерян}:
          никто не знает, какую задачу он решает и почему именно так.
  \end{itemize}
\end{frame}

% ---------------------------------------------------------------- 6
\begin{frame}{Layered Decision Records}
  \vspace{0.1cm}

  LDR собирается из трёх давно известных практик\footnote{T.\ Torres ---
  \href{https://www.producttalk.org/}{\texttt{producttalk.org}}; \ M.\ Nygard, 2011 ---
  \href{https://cognitect.com/blog/2011/11/15/documenting-architecture-decisions}{\texttt{cognitect.com}}.}:

  \vspace{0.1cm}

  \begin{columns}[t, onlytextwidth]
    \column{0.31\textwidth}
      \textbf{\small Continuous Discovery}\\
      \textcolor{mutedText}{\rule{\linewidth}{0.4pt}}\\[4pt]
      \footnotesize Teresa Torres.\\[4pt]
      Сначала фактами доказываем, что проблема есть, ---
      и только потом проектируем решение.

    \column{0.31\textwidth}
      \textbf{\small Hypothesis-Driven}\\
      \textcolor{mutedText}{\rule{\linewidth}{0.4pt}}\\[4pt]
      \footnotesize Идея входит в работу как \emph{опровержимая} гипотеза,
      а не как требование.

    \column{0.31\textwidth}
      \textbf{\small ADR}\\
      \textcolor{mutedText}{\rule{\linewidth}{0.4pt}}\\[4pt]
      \footnotesize Michael Nygard, 2011.\\[4pt]
      Принятое --- заморожено. Передумали --- пишем новое взамен,
      старое не правим.
  \end{columns}

  \vspace{0.25cm}
  \begin{center}
  \begin{minipage}{0.9\textwidth}
  \small
  Своё в LDR одно: \textbf{контекст течёт строго в одну сторону} ---
  каждый уровень добавляет только то, чего нет уровнем выше.
  Ни дублей, ни циклов.
  \end{minipage}
  \end{center}
\end{frame}

% ---------------------------------------------------------------- 7
\begin{frame}{Цикл артефактов}
  \begin{center}
  \begin{tikzpicture}
    % Single-line nodes: readable from the back row, and they leave room for
    % honest spacing between levels instead of boxes touching.
    \node[artifact accepted, minimum width=6.2cm] (prob) at (0, 2.6)
      {\artid{PROB-0001} \, thin-file applicants mispriced};

    \node[artifact accepted, minimum width=6.2cm] (met) at (0, 1.5)
      {\artid{MET-0001} \, precision @ thin-file};

    \node[artifact accepted, minimum width=3.9cm] (sol) at (-2.0, 0.4)
      {\artid{SOL-0001} \, velocity};
    \node[artifact rejected, minimum width=3.9cm] (sol2) at (2.4, 0.4)
      {\artid{SOL-0002} \, bureau};

    \node[artifact accepted, minimum width=3.9cm] (feat) at (-2.0, -0.7)
      {\artid{FEAT-0001} \, velocity v1};

    \node[artifact task, minimum width=2.0cm] (task) at (-4.3, -1.9) {\artid{TASK-0001}};
    \node[artifact task, minimum width=2.0cm] (qa)   at (-2.0, -1.9) {\artid{QA-0001}};
    \node[artifact task, minimum width=2.0cm] (rel)  at (0.3, -1.9) {\artid{REL-0001}};

    \draw[fk arrow] (met.north)  -- (prob.south);
    \draw[fk arrow] (sol.north)  -- ([xshift=-2cm]met.south);
    \draw[fk arrow] (sol2.north) -- ([xshift=2cm]met.south);
    \draw[fk arrow] (feat.north) -- (sol.south);
    \draw[fk arrow] (task.north) -- ([xshift=-1.2cm]feat.south);
    \draw[fk arrow] (qa.north)   -- (feat.south);
    \draw[fk arrow] (rel.north)  -- ([xshift=1.2cm]feat.south);
  \end{tikzpicture}
  \end{center}

  \begin{center}
  \footnotesize\textcolor{mutedText}{%
    Каждая стрелка --- внешний ключ, и идёт она \emph{от} того артефакта,
    который хранит ссылку. Корень дерева --- проблема, а не метрика.}
  \end{center}
\end{frame}

% ---------------------------------------------------------------- 8
\begin{frame}{Альтернативы не исчезают}
  \begin{center}
  \begin{tikzpicture}
    \node[artifact accepted, minimum width=9.5cm, minimum height=0.8cm] (hyp) at (0, 2.2)
      {\artid{PROB-0001}: thin-file applicants mispriced};

    \node[artifact accepted, minimum width=9cm, minimum height=0.7cm,
          align=left, anchor=west] (s1) at (-3.5, 1.0) {%
      \quad\artid{SOL-0001} transaction velocity \hfill \textbf{ACCEPTED}\quad
    };
    \node[artifact rejected, minimum width=9cm, minimum height=1.1cm,
          align=left, anchor=west] (s2) at (-3.5, -0.3) {%
      \quad\artid{SOL-0002} third-party bureau enrichment \hfill REJECTED\quad\\
      \quad\footnotesize\textit{rationale: fails GDPR review for prospect data}
    };
    \node[artifact rejected, minimum width=9cm, minimum height=1.1cm,
          align=left, anchor=west] (s3) at (-3.5, -1.7) {%
      \quad\artid{SOL-0003} graph-based network features \hfill REJECTED\quad\\
      \quad\footnotesize\textit{rationale: needs 6-month historical reload}
    };

    \draw[line width=0.7pt, black!50] (hyp.south) -- (-4, 1.8) -- (-4, -1.7);
    \draw[fk arrow] (-4, 1.0)  -- (s1.west);
    \draw[fk arrow] (-4, -0.3) -- (s2.west);
    \draw[fk arrow] (-4, -1.7) -- (s3.west);
  \end{tikzpicture}
  \end{center}

  \begin{center}
  \textcolor{mutedText}{\small Пройдёт два года, придёт новый человек и спросит:
  \textit{<<а почему тогда не взяли bureau enrichment?>>} Ответ уже лежит в
  репозитории --- вместе с причиной.}
  \end{center}
\end{frame}

% ---------------------------------------------------------------- 9
\begin{frame}{pet}
  \vspace{0.4cm}
  \begin{itemize}
    \setlength\itemsep{0.45em}
    \item \textbf{pet} --- мой инструмент для проверки этой гипотезы.
          Он сырой, и я его первый --- и пока единственный --- пользователь.
    \item Все артефакты --- обычный markdown в том же git-репозитории, что и код.
          Целостность ссылок между ними проверяет CI.
    \item Сам инструмент написан на TypeScript, а репозиторий в демо --- на \textbf{Python}.
          Это не имеет значения: артефакты лежат в \texttt{doc/},
          и язык проекта роли не играет.
  \end{itemize}
\end{frame}

% ---------------------------------------------------------------- 10
\begin{frame}{Демо: проблема и метрика}
  \shot{01-discovery-start}
\end{frame}

% ---------------------------------------------------------------- 11
\begin{frame}{Демо: агент заполняет Evidence}
  \shot{02-researcher}
\end{frame}

% ---------------------------------------------------------------- 12
\begin{frame}{Демо: решает человек, а не агент}
  \shot{04-reject-alternative}
\end{frame}

% ---------------------------------------------------------------- 13
\begin{frame}{Демо: отвергнутая альтернатива остаётся}
  \shot{05-alternatives-preserved}
\end{frame}

% ---------------------------------------------------------------- 14
\begin{frame}{Демо: что легло в git}
  \shot{06-pipeline}
\end{frame}

% ---------------------------------------------------------------- 15
\begin{frame}{Демо: принятое не переписывается}
  \shot{08-immutability}
\end{frame}

% ---------------------------------------------------------------- 16
\begin{frame}{Чем это отличается от ADR и spec-driven}
  \vspace{0.25cm}
  \begin{columns}[t, onlytextwidth]

    \column{0.47\textwidth}
      \textbf{\small ADR}\\
      \textcolor{mutedText}{\rule{\linewidth}{0.4pt}}\\[5pt]
      \footnotesize Решает ровно половину задачи: фиксирует \emph{архитектурные}
      решения --- и делает это хорошо.\\[5pt]
      Но у него нет связи с продуктовым <<зачем>>: ни метрики, ни гипотезы,
      ни уровня выше.

    \column{0.47\textwidth}
      \textbf{\small Spec-driven: Kiro, Spec Kit}\\
      \textcolor{mutedText}{\rule{\linewidth}{0.4pt}}\\[5pt]
      \footnotesize Спека \emph{запускает} работу --- и на этом её жизнь
      заканчивается: дальше она протухает, и источником правды снова
      становится код.\\[5pt]
      Долгоживущие спеки сегодня пробует только Tessl.

  \end{columns}

  \vspace{0.4cm}
  \small Мейнстримный spec-driven development копит ровно тот долг, о котором
  этот доклад: спецификация там --- одноразовый запуск, а не запись
  решения\footnote{Разбор трёх SDD-инструментов: B.\ Böckeler.
  \textit{Understanding Spec-Driven-Development: Kiro, spec-kit, and Tessl.}
  martinfowler.com, 2025.
  \href{https://martinfowler.com/articles/exploring-gen-ai/sdd-3-tools.html}{\texttt{martinfowler.com/articles/exploring-gen-ai/sdd-3-tools.html}}}.
\end{frame}

% ---------------------------------------------------------------- 17
\begin{frame}{Возражения}
  \vspace{0.25cm}
  \begin{itemize}
    \setlength\itemsep{0.5em}
    \small
    \item \textbf{<<Это просто провал инженерной дисциплины>>}\\
          \footnotesize Возможно. Но дисциплина не масштабируется и не переживает
          смену команды. Нужен механизм, который работает без героев, ---
          а не благие намерения.

    \item \textbf{<<AI сам восстановит rationale по коду>>}\\
          \footnotesize Восстановит --- но только \emph{что} сделано, не \emph{почему}.
          А отвергнутой альтернативы в коде нет вообще: её невозможно вывести
          из того, что осталось.

    \item \textbf{<<Resist the automation of understanding>>}\footnote{М.-Э.\ Стори
          об этом же: \textit{What I'm Hearing About Cognitive Debt (So Far)},
          2026. \href{https://margaretstorey.com/blog/2026/02/18/cognitive-debt-revisited/}{\texttt{margaretstorey.com}}}\\
          \footnotesize Это возражение Стори бьёт прямо по pet: артефакты пишет агент.
          Мой ответ --- человек жмёт \texttt{accept} на каждом шаге. Агент готовит,
          решает человек; без этого запись превращается в ритуал.
  \end{itemize}
\end{frame}

% ---------------------------------------------------------------- 18
\begin{frame}{Честная граница}
  \vspace{0.3cm}
  \begin{itemize}
    \setlength\itemsep{0.45em}
    \item \textbf{Discovery действительно догфудится:} проблема, метрики,
          решения и фичи этого репозитория заведены через сам инструмент.
    \item \textbf{Delivery --- пока нет:} задачи, QA-планы и релизы
          построены мимо пайплайна.
    \item Это проверяется одной командой \texttt{ls} --- поэтому я говорю
          об этом сам, не дожидаясь вопроса из зала.
  \end{itemize}

  \vspace{0.5cm}
  \begin{center}
  \begin{minipage}{0.9\textwidth}
  \footnotesize\textcolor{mutedText}{%
  Инструмент против intent debt обязан выдерживать собственную проверку.
  Там, где он её пока не выдерживает, честнее сказать это со сцены раньше,
  чем спросят из зала.}
  \end{minipage}
  \end{center}
\end{frame}

% ---------------------------------------------------------------- 19
\begin{frame}{Что померить и что попробовать}
  \vspace{0.25cm}

  \textbf{\small Метрики intent debt в вашей команде}\\[-6pt]
  \textcolor{mutedText}{\rule{\linewidth}{0.4pt}}\\[3pt]
  \begin{itemize}
    \setlength\itemsep{0.15em}
    \footnotesize
    \item У какой доли решений записана причина --- а не только результат?
    \item Сколько времени занимает ответ на вопрос <<почему так, а не иначе?>>
    \item Сколько отвергнутых альтернатив за прошлый квартал вы можете вспомнить?
  \end{itemize}

  \vspace{0.4cm}
  \textbf{\small Что попробовать уже завтра}\\[-6pt]
  \textcolor{mutedText}{\rule{\linewidth}{0.4pt}}\\[3pt]
  \begin{itemize}
    \setlength\itemsep{0.15em}
    \footnotesize
    \item Написать один ADR на ближайшее архитектурное решение. Один файл, пять минут.
    \item Записать не только выбранный вариант, но и отвергнутый --- вместе с причиной.
    \item Проверка простая: сможет ли новый человек ответить <<почему не X?>>,
          не приходя к вам.
  \end{itemize}
\end{frame}

% ---------------------------------------------------------------- 20
\begin{frame}{Потрогать руками}
  \begin{center}
    {\hypersetup{urlcolor=black}%
     \qrcode[height=3.7cm]{https://github.com/smirnoffmg/pet-example}}\\[0.4cm]
    \href{https://github.com/smirnoffmg/pet-example}{\texttt{github.com/smirnoffmg/pet-example}}\\[0.15cm]
    {\small\textcolor{mutedText}{демо-репозиторий доклада: Python-сервис и весь пайплайн
    артефактов в \texttt{doc/}}}\\[0.45cm]
    {\small инструмент: \href{https://github.com/smirnoffmg/pet}{\texttt{github.com/smirnoffmg/pet}}}
  \end{center}
\end{frame}

% ---------------------------------------------------------------- 21
\begin{frame}{Источники}
  \scriptsize
  \begin{itemize}
    \setlength\itemsep{0.2em}

    \item M.-A.\ Storey. \textit{From Technical Debt to Cognitive and Intent
          Debt: Rethinking Software Health in the Age of AI.}
          ACM Queue, 2026.
          \href{https://doi.org/10.1145/3807966}{\texttt{doi.org/10.1145/3807966}}
          \textcolor{mutedText}{(препринт:
          \href{https://arxiv.org/abs/2603.22106}{\texttt{arXiv:2603.22106}})}

    \item M.-A.\ Storey. \textit{What I'm Hearing About Cognitive Debt (So Far).}
          2026.
          \href{https://margaretstorey.com/blog/2026/02/18/cognitive-debt-revisited/}{\texttt{margaretstorey.com/blog/2026/02/18/cognitive-debt-revisited}}

    \item P.\ Naur. \textit{Programming as Theory Building.}
          Microprocessing and Microprogramming, 1985.
          \href{https://pages.cs.wisc.edu/~remzi/Naur.pdf}{\texttt{pages.cs.wisc.edu/\textasciitilde remzi/Naur.pdf}}

    \item M.\ Nygard. \textit{Documenting Architecture Decisions.} 2011.
          \href{https://cognitect.com/blog/2011/11/15/documenting-architecture-decisions}{\texttt{cognitect.com/blog/2011/11/15/documenting-architecture-decisions}}

    \item T.\ Torres. \textit{Continuous Discovery Habits.} 2021.
          \href{https://www.producttalk.org/}{\texttt{producttalk.org}}

    \item H.\ Cummins. \textit{Intent Debt in Software Development}
          \textcolor{mutedText}{(пост в LinkedIn, 2026)}. Её кейноут на QCon London 2026:
          \textit{The Free-Lunch Guide to Idea Circularity} ---
          \href{https://speakerdeck.com/hollycummins/the-free-lunch-guide-to-idea-circularity}{\texttt{speakerdeck.com/hollycummins}}

    \item A.\ Osmani. \textit{The Intent Debt.}
          \href{https://addyosmani.com/blog/intent-debt/}{\texttt{addyosmani.com/blog/intent-debt}}

    \item B.\ Böckeler. \textit{Understanding Spec-Driven-Development: Kiro,
          spec-kit, and Tessl.} martinfowler.com, 2025.
          \href{https://martinfowler.com/articles/exploring-gen-ai/sdd-3-tools.html}{\texttt{martinfowler.com/articles/exploring-gen-ai/sdd-3-tools.html}}

    \item \textit{Faster Code, Deeper Debt? A Multivocal Literature Review on
          Technical Debt in LLM-Assisted Software Development.}
          \href{https://arxiv.org/abs/2606.14796}{\texttt{arXiv:2606.14796}}

  \end{itemize}
\end{frame}

% ---------------------------------------------------------------- 22
\begin{frame}[plain]
  \vfill
  \begin{center}
    {\Huge\textbf{Спасибо!}}\\[0.7cm]
    {\large Максим Смирнов}\\[0.5cm]
    {\hypersetup{urlcolor=black}%
     \qrcode[height=3.2cm]{https://smirnoffmg.dev}}\\[0.35cm]
    \href{https://smirnoffmg.dev}{\texttt{\large smirnoffmg.dev}}
  \end{center}
  \vfill
\end{frame}

\end{document}
